\chapter{Практика}

\section{Задача}
С помощью Wireshark изучить пакеты DHCP

\section{DHCP Discover}

Клиент обращается к DHCP серверу и просит IP, DNS, маску и т.д

\begin{figure}[H] 
    \centering
    \includegraphics[width=1.0\textwidth]{dis.png}
    \caption{DHCP Discover}
\end{figure}

Основные поля:

\begin{enumerate}
    \item Client IP address - ip адрес клиента, но т.к его нет, то отправляется просто 0.0.0.0
    \item Client MAC address - MAC-адрес клиента. Отправляется, чтобы DHCP сервер знал, кому отправить offer
    \item Requested IP address - ip, который хотел бы получить клиент.
\end{enumerate}

\section{DHCP Offer}

DHCP сервер отвечает клиенту и отправляет IP, который он может предложить.

\begin{figure}[H] 
    \centering
    \includegraphics[width=1.0\textwidth]{of.png}
    \caption{DHCP Offer}
\end{figure}

Основные поля:

\begin{enumerate}
    \item Your IP address - адрес, который предлагает DHCP сервер (тот же, что и просили)
    \item DHCP server ID - адрес DHCP сервера
    \item Address Lease Time - время аренды адреса
    \item Renewal Time Value - время, через которое будет требоваться продление. Обозначается как $T_1$, обычно равно 50\% от времени аренды.
    \item Rebinding Time Value - время, когда повторно будет происходить продление, если в первый раз не получилоь продлить время. Обозначается как $T_2$, обычно равно 87.5\% от времени аренды.
\end{enumerate}

\section{DHCP Request}

Клиент отвечает DHCP серверу на предложение (берет IP или нет)

\begin{figure}[H] 
    \centering
    \includegraphics[width=1.0\textwidth]{req.png}
    \caption{DHCP Request}
\end{figure}

Основные поля:

\begin{enumerate}
    \item Requested IP address - адрес, который клиент соглашается взять.
\end{enumerate}

\section{DHCP ACK}

DHCP сервер подтверждает передачу IP клиентскому устройству

\begin{figure}[H] 
    \centering
    \includegraphics[width=1.0\textwidth]{ack.png}
    \caption{DHCP Request}
\end{figure}

\endinput